% =============================================================================
%  2-Introduccion.tex — Neurodiversidad, la escuela como entorno sensorial
%  y la inclusión como eliminación de estímulos que limitan el aprendizaje
% =============================================================================
%  NOTA: citas en traducción propia; el original en inglés va en comentario
%        (% ORIG:) junto a cada cita, para verificación.
%  EDITAR: Pellicano y den Houting (2022) — falta número de página exacto de
%          cada cita (se consultó la versión HTML/XML de PMC, sin paginación).
% =============================================================================

\section{Introducción}

Comprender la neurodiversidad es una tarea que atañe tanto a quien enseña como a quien diseña el espacio. \citet{pellicano2022} describen el paradigma de la neurodiversidad como aquel que \enquote{rechaza la idea de que la divergencia respecto a la norma sea un defecto que requiera corrección}\footnote{Todas las citas de este documento son traducción propia del original en inglés.}.
% ORIG: rejects the view that divergence from the norm is a flaw requiring correction
Desde esta perspectiva, los autores sostienen que \enquote{cualquier \enquote{discapacidad} no se entiende mejor como resultado de las características únicas de un individuo, sino como resultado de un entorno que no acomoda eficazmente esas características}, y advierten:
% ORIG: any ‘disability’ is not best understood as the result of an individual’s
%       unique characteristics, but rather as a result of an environment that
%       does not effectively accommodate those characteristics

\begin{displayquote}
Los entornos físicos y sociales en los que todos vivimos están, en general, diseñados para satisfacer las necesidades de quienes se ubican dentro del rango típico del neurodesarrollo; esos mismos entornos son con frecuencia subóptimos, e incluso hostiles, para las personas neurodivergentes, y necesitan ajustarse si se quiere que estas personas lleven una buena vida \citep{pellicano2022}.
\end{displayquote}
% ORIG: The physical and social environments within which we all live are
%       generally designed to meet the needs of those who fall within the
%       typical range of neurodevelopment; these same environments are often
%       suboptimal, and even hostile, for neurodivergent people and need to be
%       adjusted if such people are to lead good lives

Esta lectura coloca al entorno construido en el centro de la discusión. Desde la arquitectura, \citet[p.~143]{mostafa2014} lo plantea así: \enquote{La arquitectura, como ciencia, se ocupa de la manipulación del entorno físico para facilitar ciertas funciones y propiciar comportamientos deseados. Este entorno se compone principalmente de elementos sensoriales: texturas, colores, patrones, acústica, etc.}
% ORIG: Architecture, as a science, deals with the manipulation of the physical
%       environment to facilitate certain functions and elicit intended
%       behavior. This environment is comprised primarily of sensory elements-
%       textures, colors, patterns, acoustics etc.

La escuela es uno de los entornos donde este desajuste se vuelve más evidente. \citet{pellicano2022} la describen en estos términos: \enquote{Los entornos escolares tradicionales y convencionales son un ejemplo de este tipo de entorno: suelen ser físicamente grandes, ruidosos y caóticos, exigen transiciones frecuentes dentro de las clases y entre ellas a lo largo de la jornada escolar, e implican una serie de reglas y expectativas sociales implícitas}.
% ORIG: Traditional, mainstream school settings are an example of one such
%       environment: they are often physically large, noisy and chaotic,
%       require frequent transitions within and between classes throughout the
%       school day, and entail a host of implicit social rules and expectations
Añaden que \enquote{los entornos físicos, como las escuelas, diseñados pensando en las necesidades neurotípicas pueden dejar a las personas autistas en un estado de incomodidad y angustia sensorial, y los entornos sociales regidos por reglas neurotípicas suelen ser igualmente inaccesibles}.
% ORIG: Physical environments, like schools, that are designed with neurotypical
%       needs in mind can leave autistic people in sensory discomfort and
%       distress, and social environments governed by neurotypical rules are
%       often similarly inaccessible

\citet[p.~2]{jones2020} ilustran el mismo fenómeno: \enquote{un niño con hipersensibilidad puede funcionar plenamente en casa, pero experimentar \enquote{discapacidad} al ser colocado en un aula concurrida con estímulos sensoriales significativos, como luces brillantes y niños ruidosos}.
% ORIG: a child with hypersensitivity may be able to fully function at home but
%       may experience ‘disability’ when placed in a busy classroom with
%       significant sensory inputs such as bright lights and noisy children
En su estudio, un testimonio docente lo resume así: \enquote{Con frecuencia los entornos del aula están demasiado saturados visualmente, lo que significa que muchos jóvenes no saben a qué prestar atención y se sienten abrumados} \citep[p.~6]{jones2020}.
% ORIG: Often classroom environments are too visually busy which means many
%       youngsters don't know what to attend to and are overwhelmed
El efecto de la carga visual se ha documentado incluso en población infantil en general: \enquote{Los niños se distrajeron más con el entorno visual, pasaron más tiempo fuera de la tarea y mostraron menores avances en el aprendizaje cuando las paredes estaban muy decoradas que cuando se retiraron las decoraciones} \citep[p.~1362]{fisher2014}.
% ORIG: Children were more distracted by the visual environment, spent more time
%       off task, and demonstrated smaller learning gains when the walls were
%       highly decorated than when the decorations were removed
En el plano auditivo, \citet[p.~564]{ashburner2008} concluyen que \enquote{los niños que tienen dificultad para procesar instrucciones verbales en entornos ruidosos y que a menudo se concentran en conductas de búsqueda sensorial parecen tener mayor probabilidad de un bajo rendimiento académico}.
% ORIG: Children who have difficulty processing verbal instructions in noisy
%       environments and who often focus on sensory-seeking behaviors appear
%       more likely to underachieve academically

Si la barrera está en el entorno, la inclusión empieza por identificarla. Al revisar los factores que modulan el impacto sensorial en el aprendizaje, \citet[p.~11]{jones2020} reportan \enquote{la agencia, la previsibilidad, el diseño del aula, los recursos escolares y el conocimiento del personal como elementos clave para minimizar el impacto de las diferencias sensoriales}.
% ORIG: agency, predictability, classroom design, school resources and staff
%       knowledge as key in minimizing the impact of sensory differences
Esta guía parte de esa premisa: antes de pedir a quien aprende que se adapte, hay que reconocer qué estímulos del aula limitan su aprendizaje. Para ello, la sesión convierte el aula en un entorno inmersivo que permite experimentar, de forma breve y controlada, algunas de esas barreras invisibles.
